\subsubsection{Measuring \texorpdfstring{$\ket{+}$}{|+>} with the Pauli-Z Observable}

After introducing observables, projectors, probabilities, and expectation values, it is useful to analyze a concrete example that combines all these concepts.

\highspace
Consider the quantum state:
\begin{equation*}
    \ket{+}
    =
    \frac{1}{\sqrt{2}}
    \left(
        \ket{0}
        +
        \ket{1}
    \right)
\end{equation*}
And the Pauli-Z observable:
\begin{equation*}
    Z
    =
    \begin{bmatrix}
        1 & 0 \\
        0 & -1
    \end{bmatrix}
\end{equation*}
\hl{Calculate the possible measurement outcomes, their probabilities, the post-measurement states, and the expectation value of $Z$ in the state $\ket{+}$.}

\highspace
From the spectral decomposition (\autopageref{eq:spectral-decomposition}) of $Z$, we know that:
\begin{itemize}
    \item The \textbf{eigenvalues} of $Z$ are $+1$ and $-1$:
    \begin{align*}
        \det\left(Z - \lambda I\right) &= 0 \\
        \det\left(\begin{bmatrix}
            1 - \lambda & 0 \\
            0 & -1 - \lambda
        \end{bmatrix}\right) &= 0 \\
        (1 - \lambda)(-1 - \lambda) &= 0 \\
        \lambda &= \pm 1
    \end{align*}
    
    \item The \textbf{first eigenvector} \ket{v_1} is \ket{0}:
    \begin{align*}
        \left(Z - 1 I\right)\ket{v_1} &= 0 \\
        \begin{bmatrix}
            1 - 1 & 0 \\
            0 & -1 - 1
        \end{bmatrix}
        \begin{bmatrix}
            x \\
            y
        \end{bmatrix} &= 0 \\
        \begin{bmatrix}
            0 & 0 \\
            0 & -2
        \end{bmatrix}
        \begin{bmatrix}
            x \\
            y
        \end{bmatrix} &= 0 \\
        \begin{cases}
            0 \cdot x + 0 \cdot y = 0 \\
            0 \cdot x - 2 \cdot y = 0
        \end{cases} &\implies
        y = 0
    \end{align*}
    So any vector of the form $\begin{pmatrix}
        x \\
        y
    \end{pmatrix} \xrightarrow{y=0} \begin{bmatrix}
        x \\
        0
    \end{bmatrix}$ is an eigenvector. We choose the simplest one $\left(1, 0\right)$, and then we normalize it:
    \begin{equation*}
        \left\|\begin{bmatrix}
            1 \\
            0
        \end{bmatrix}\right\| = \sqrt{1^2 + 0^2} = 1
    \end{equation*}
    Therefore, the normalized eigenvector is:
    \begin{equation*}
        \ket{v_1} = 1 \cdot \begin{bmatrix}
            1 \\
            0
        \end{bmatrix} = \ket{0}
    \end{equation*}

    \item Similarly, the \textbf{second eigenvector} \ket{v_2} is \ket{1}:
    \begin{align*}
        \left(Z + 1 I\right)\ket{v_2} &= 0 \\
        \begin{bmatrix}
            1 + 1 & 0 \\
            0 & -1 + 1
        \end{bmatrix}
        \begin{bmatrix}
            x \\
            y
        \end{bmatrix} &= 0 \\
        \begin{bmatrix}
            2 & 0 \\
            0 & 0
        \end{bmatrix}
        \begin{bmatrix}
            x \\
            y
        \end{bmatrix} &= 0 \\
        \begin{cases}
            2 \cdot x + 0 \cdot y = 0 \\
            0 \cdot x + 0 \cdot y = 0
        \end{cases} &\implies
        x = 0
    \end{align*}
    So any vector of the form $\begin{pmatrix}
        x \\
        y
    \end{pmatrix} \xrightarrow{x=0} \begin{bmatrix}
        0 \\
        y
    \end{bmatrix}$ is an eigenvector. We choose the simplest one $\left(0, 1\right)$, and then we normalize it:
    \begin{equation*}
        \left\|\begin{bmatrix}
            0 \\
            1
        \end{bmatrix}\right\| = \sqrt{0^2 + 1^2} = 1
    \end{equation*}
    Therefore, the normalized eigenvector is:
    \begin{equation*}
        \ket{v_2} = 1 \cdot \begin{bmatrix}
            0 \\
            1
        \end{bmatrix} = \ket{1}
    \end{equation*}
\end{itemize}

Therefore:
\begin{itemize}
    \item The possible \hl{measurement outcomes} are \hl{$+1$} and \hl{$-1$};
    \item The corresponding \hl{eigenstates} are \hl{$\ket{0}$} and \hl{$\ket{1}$};
    \item The \hl{measurement operators} are:
    \begin{equation*}
        M_0 = \ket{0}\bra{0},
        \qquad
        M_1 = \ket{1}\bra{1}
    \end{equation*}
\end{itemize}

\begin{flushleft}
    \textcolor{Green3}{\faIcon{book} \textbf{Measurement Probabilities}}
\end{flushleft}
The state $\ket{+}$ is an equal superposition of $\ket{0}$ and $\ket{1}$:
\begin{equation*}
    \ket{+}
    =
    \frac{1}{\sqrt{2}}\ket{0}
    +
    \frac{1}{\sqrt{2}}\ket{1}
\end{equation*}
Therefore, the amplitudes are:
\begin{equation*}
    \alpha = \frac{1}{\sqrt{2}},
    \qquad
    \beta = \frac{1}{\sqrt{2}}
\end{equation*}
Using the Born rule:
\begin{equation*}
    p_0 = |\alpha|^2 = \frac{1}{2},
    \qquad
    p_1 = |\beta|^2 = \frac{1}{2}
\end{equation*}
Thus:
\begin{itemize}
    \item There is a 50\% probability of obtaining the outcome $+1$;
    \item There is a 50\% probability of obtaining the outcome $-1$.
\end{itemize}

\begin{flushleft}
    \textcolor{Green3}{\faIcon{book} \textbf{Post-Measurement States}}
\end{flushleft}
If the measurement outcome is:
\begin{itemize}
    \item $+1$, the state collapses to $\ket{0}$;
    \item $-1$, the state collapses to $\ket{1}$.
\end{itemize}

\highspace
\begin{flushleft}
    \textcolor{Green3}{\faIcon{book} \textbf{Expectation Value}}
\end{flushleft}
The expectation value of $Z$ in the state $\ket{+}$ is:
\begin{equation*}
    \bra{+} Z \ket{+}
    =
    (+1)\cdot \frac{1}{2}
    +
    (-1)\cdot \frac{1}{2}
    =
    0
\end{equation*}

\begin{flushleft}
    \textcolor{Red2}{\faIcon{exclamation-triangle} \textbf{Remark}}
\end{flushleft}
The expectation value is $0$, but this does \textbf{not} mean that the measurement can return the value $0$. The only possible outcomes are:
\begin{equation*}
    +1
    \quad
    \text{and}
    \quad
    -1
\end{equation*}
The value $0$ is simply the probability-weighted average of these outcomes over many repeated measurements.

\highspace
\begin{flushleft}
    \textcolor{Green3}{\faIcon{question-circle} \textbf{What if the state to be measured is already an eigenstate of the observable?}}
\end{flushleft}
Let's assume that now we want to calculate the measurement outcomes, probabilities, post-measurement states, and expectation value when the input state is $\ket{0}$ or $\ket{1}$ instead of $\ket{+}$. How we have seen, $\ket{0}$ and $\ket{1}$ are eigenvectors of $Z$ with eigenvalues $+1$ and $-1$, respectively. This means that they are already measurement states of the observable $Z$. Let's analyze what happens in these cases.
\begin{itemize}
    \item \textbf{Measuring \ket{0}}: the measurement probabilities of \ket{0} are:
    \begin{equation*}
        \ket{0} = 1 \cdot \ket{0} + 0 \cdot \ket{1}
    \end{equation*}
    Therefore:
    \begin{equation*}
        p_0 = \left|\alpha\right|^2 = 1,
        \qquad
        p_1 = \left|\beta\right|^2 = 0
    \end{equation*}
    Using observable $Z$, there is a 100\% probability of obtaining the outcome $+1$ and a 0\% probability of obtaining the outcome $-1$. Moreover, the (post-)measurement outcome is always $+1$, and the post-measurement state remains $\ket{0}$ because we have 100\% probability of collapsing to $\ket{0}$ and 0\% probability of collapsing to $\ket{1}$. Finally, the expectation value is:
    \begin{equation*}
        \bra{0} Z \ket{0} = \left(+1\right) \cdot 1 + \left(-1\right) \cdot 0 = +1
    \end{equation*}
    And indeed, the expectation value confirms that the measurement outcome is always $+1$.
    
    \item \textbf{Measuring \ket{1}}: the measurement probabilities of \ket{1} are:
    \begin{equation*}
        \ket{1} = 0 \cdot \ket{0} + 1 \cdot \ket{1}
    \end{equation*}
    Therefore:
    \begin{equation*}
        p_0 = \left|\alpha\right|^2 = 0,
        \qquad
        p_1 = \left|\beta\right|^2 = 1
    \end{equation*}
    Using observable $Z$, there is a 0\% probability of obtaining the outcome $+1$ and a 100\% probability of obtaining the outcome $-1$. Moreover, the (post-)measurement outcome is always $-1$, and the post-measurement state remains $\ket{1}$ because we have 0\% probability of collapsing to $\ket{0}$ and 100\% probability of collapsing to $\ket{1}$. Finally, the expectation value is:
    \begin{equation*}
        \bra{1} Z \ket{1} = \left(+1\right) \cdot 0 + \left(-1\right) \cdot 1 = -1
    \end{equation*}
    And indeed, the expectation value confirms that the measurement outcome is always $-1$.
\end{itemize}
In summary, \hl{when a quantum state is already an eigenvector of the observable being measured}:
\begin{itemize}
    \item The \textbf{measurement} outcome is \textbf{deterministic};
    \item The \textbf{probability} of the corresponding eigenvalue is \textbf{1};
    \item The \textbf{state does not change} after the measurement;
    \item The \textbf{expectation value} is exactly the associated \textbf{eigenvalue}.
\end{itemize}
In contrast, when the state is a superposition of eigenvectors (as in the case of $\ket{+}$), the measurement outcome is probabilistic and the expectation value becomes a weighted average of the possible outcomes.